% Part: first-order-logic
% Chapter: axiomatic-deduction
% Section: provability-quantifiers
% Verification of properties of provability needed for maximally
% consistent sets in the completeness chapter.

\documentclass[../../../include/open-logic-section]{subfiles}

\begin{document}

\olfileid{fol}{axd}{qpr}

\olsection{可推导性与量词}

\begin{explain}
  完备性定理还要求，公理演绎能够得出本节所建立的关于可证关系~$\Proves$ 的事实。
\end{explain}

\begin{thm}
\ollabel{thm:strong-generalization} 如果 $c$ 是一个不在 $\Gamma$ 或 $!A(x)$ 中出现的常项符号，并且 $\Gamma \Proves !A(c)$，那么 $\Gamma
\Proves \lforall[x][!A(x)]$。
\end{thm}

\begin{proof}
根据演绎定理，有 $\Gamma \Proves \ltrue \lif !A(c)$。由于 $c$ 不在 $\Gamma$ 或~$\top$ 中出现，我们得到 $\Gamma \Proves
\ltrue \lif \lforall[x][!A(x)]$。再次应用演绎定理，即得 $\Gamma \Proves
\lforall[x][!A(x)]$。
\end{proof}

\begin{prop}
\ollabel{prop:provability-quantifiers}
\begin{tagenumerate}{prvEx,prvAll}
\tagitem{prvEx}{$!A(t) \Proves \lexists[x][!A(x)]$.}{}

\tagitem{prvAll}{$\lforall[x][!A(x)] \Proves !A(t)$.}{}
\end{tagenumerate}
\end{prop}

\begin{proof}
\begin{tagenumerate}{prvEx,prvAll}
\tagitem{prvEx}{由 \olref[qua]{ax:q2} 及演绎定理可得。}{}
\tagitem{prvAll}{由 \olref[qua]{ax:q1} 及演绎定理可得。}{}
\end{tagenumerate}
\end{proof}

\end{document}
